\subsection{Gaussian ray weights and their progression main terms}
\label{subsec:gaussian-progression-mainterms}

Let $E=\Q(i)$ and let $H$ be a fixed positive integer divisible by $16$.
Let $K/E$ be the ray class field for $H\Z[i]$.  This field is Galois over
$\Q$, since complex conjugation preserves both the modulus and its
principal ray subgroup.  Put
\[
 A=\operatorname{Gal}(K/E)
   \simeq (\Z[i]/H\Z[i])^\times/\{1,-1,i,-i\},
 \qquad G=\operatorname{Gal}(K/\Q).
\]
We use arithmetic Frobenius throughout.  Let $C\subseteq A$ be invariant
under complex conjugation, so that $C$ is a union of conjugacy classes in
$G$.  Define $\Lambda'_C(p)=\log p$ for unramified rational primes whose
Frobenius lies in $C$, and put $\Lambda'_C(n)=0$ otherwise.  In particular
this function is supported on rational primes which split in $E$.

For ideals coprime to $H$ let $[\mathfrak a]\in A$ denote their ray class,
and set
\[
 T_C(n)=\sum_{\substack{\mathrm N\mathfrak a=n\\(\mathfrak a,H)=1}}
       \Lambda_E(\mathfrak a)1_C([\mathfrak a]).
\]
Here $\Lambda_E(\mathfrak p^j)=\log\mathrm N\mathfrak p$ and it is zero
on all other ideals.  Every unramified split rational prime has two
conjugate prime ideals of norm $p$, so
\begin{equation}\label{eq:gaussian-ray-prime-factor}
 T_C(p)=2\Lambda'_C(p)\qquad(p\nmid2H).
\end{equation}
The remaining difference between $T_C$ and $2\Lambda'_C$ is supported on
rational prime powers of exponent at least two and the finitely many
excluded primes.  Its total absolute value up to $X$ is
$O_H(\sqrt X\log^2(2X))$.  This bound follows by summing the at most two
prime ideals above each rational prime and their powers; no density
theorem is used in this comparison.

For $\chi\in\widehat A$ put
\[
 c_C(\chi)=|A|^{-1}\sum_{a\in C}\overline{\chi(a)},\qquad
 T_\chi(n)=\sum_{\mathrm N\mathfrak a=n}
                  \Lambda_E(\mathfrak a)\chi(\mathfrak a),
\]
where each character is extended by zero on ideals not coprime to $H$.
Finite Fourier inversion gives the exact identity
\begin{equation}\label{eq:gaussian-ray-fourier}
                 T_C(n)=\sum_{\chi\in\widehat A}c_C(\chi)T_\chi(n).
\end{equation}
The field and all these characters are fixed before $X$ tends to infinity.

\begin{lemma}[Fixed-field Hecke Siegel--Walfisz estimate]
\label{lem:gaussian-hecke-sw}
Let $E=\Q(i)$, let $\chi$ be a fixed finite-order ray character of $E$
with finite modulus $\mathfrak m$, and let $B>0$.
There is $c=c(\chi,B)>0$ such that, uniformly for $X\geq3$,
$1\leq Y\leq X$, positive integers $q\leq(\log X)^B$, and Dirichlet
characters $\rho$ modulo $q$,
\begin{equation}\label{eq:gaussian-hecke-sw-estimate}
 \sum_{\mathrm N\mathfrak a\leq Y}
   \Lambda_E(\mathfrak a)\chi(\mathfrak a)\rho(\mathrm N\mathfrak a)
 =r_{\chi,\rho}Y+
 O_{\chi,B}\bigl(X\exp(-c\sqrt{\log X})\bigr).
\end{equation}
Here $\Lambda_E(\mathfrak p^j)=\log\mathrm N\mathfrak p$ for a prime
ideal $\mathfrak p$ and $j\geq1$, and it is zero otherwise; characters
are extended by zero on ideals not coprime to their moduli.
The number $r_{\chi,\rho}$ is one if the primitive Hecke character
inducing $\chi(\rho\circ\mathrm N)$ is trivial, and zero otherwise.
There is also $c_B>0$ such that, with the same ranges of $X,Y,q,\rho$,
\begin{equation}\label{eq:rational-sw-estimate}
 \sum_{n\leq Y}\Lambda(n)\rho(n)
 =1_{\rho\ \mathrm{principal}}Y+
 O_B\bigl(X\exp(-c_B\sqrt{\log X})\bigr).
\end{equation}
In particular, either error is $O(X(\log X)^{-D})$ for every fixed
$D>0$, with the corresponding dependence on $D$ allowed.
\end{lemma}

\begin{proof}
We use the published fixed-field estimate of Mitsui
\cite[\S1, Lemma~15, p.~12]{Mitsui1956}.
In his notation, for a fixed number field $K$, a fixed $A>0$, and any
ray character $\vartheta$ whose finite modulus has norm at most
$(\log x)^A$, it states
\begin{equation}\label{eq:mitsui-weighted-primes}
 \sum_{\mathrm N\mathfrak p\leq x}
       \vartheta(\mathfrak p)\log\mathrm N\mathfrak p
 =E(\vartheta)x+
 O_{K,A}\bigl(x\exp(-c_{K,A}\sqrt{\log x})\bigr),
\end{equation}
where $E(\vartheta)=1$ for the principal character and is zero otherwise.
The sum is over prime ideals of every residue degree.
Mitsui defines these characters and $E(\vartheta)$ on p.~2, includes
imprimitive characters, and states the uniformity in the character
and its modulus immediately before Lemma~14 on p.~11.
Infinite primes may occur in the ray modulus; only its finite part
enters the norm bound.

For completeness, the elementary passage from Mitsui's proved
smoothed estimate (Lemma~14, pp.~11--12) to
\eqref{eq:mitsui-weighted-primes} also applies to complex characters.
Put
\[
 P(x,\vartheta)=\sum_{\mathrm N\mathfrak p\leq x}
  \vartheta(\mathfrak p)\log\mathrm N\mathfrak p
                         \log\frac{x}{\mathrm N\mathfrak p}.
\]
Lemma~14 gives $P(x,\vartheta)=E(\vartheta)x+
O_{K,A}(xe^{-c_1\sqrt{\log x}})$.
Subtract its values at $xe^h$ and $x$, where
$h=e^{-c_1\sqrt{\log x}/2}$, and divide by $h$.
The resulting left side is the prime sum in
\eqref{eq:mitsui-weighted-primes} plus boundary terms of absolute
value at most
\[
 \sum_{x<\mathrm N\mathfrak p\leq xe^h}\log\mathrm N\mathfrak p
       \ll_K(xh+1)\log x.
\]
Indeed, for each integer norm there are at most $[K:\Q]$ prime ideals.
The main term is $E(\vartheta)x(e^h-1)/h=E(\vartheta)x+O(xh)$,
and the error divided by $h$ is
$O_{K,A}(xe^{-c_1\sqrt{\log x}/2})$.
Reducing the positive exponent constant proves
\eqref{eq:mitsui-weighted-primes}.

Prime powers of exponent at least two contribute
$O_K(\sqrt x\log^2(2x))$: sum over exponents $j\geq2$ and use the
elementary bound $O_K(x^{1/j})$ for the number of prime ideals of norm
at most $x^{1/j}$.
Consequently \eqref{eq:mitsui-weighted-primes} holds with the prime
sum replaced by
$\sum_{\mathrm N\mathfrak a\leq x}
\Lambda_K(\mathfrak a)\vartheta(\mathfrak a)$.

The character $\rho\circ\mathrm N$ on $E$ has ray modulus $(q)$:
if $\alpha\equiv1\pmod{q\Z[i]}$, then
$\mathrm N\alpha\equiv1\pmod q$.
Thus $\chi(\rho\circ\mathrm N)$ has a modulus dividing
$\mathfrak m(q)$, whose norm is $\mathrm N\mathfrak m\,q^2$.
For $\sqrt X\leq Y\leq X$, the bounds
$q\leq(\log X)^B$ and $\log Y\geq\frac12\log X$ imply
\[
 \mathrm N\mathfrak m\,q^2\leq(\log Y)^{2B+1}
\]
for all sufficiently large $X$, with the threshold depending only
on $\mathfrak m$ and $B$.
Apply the preceding estimate with $K=E$, $x=Y$, and $A=2B+1$.
Its error is at most the error in
\eqref{eq:gaussian-hecke-sw-estimate}, after reducing $c$.
For $1\leq Y<\sqrt X$, the elementary bound
$\sum_{\mathrm N\mathfrak a\leq Y}\Lambda_E(\mathfrak a)
\ll_E Y\log^2(2Y)$ suffices, and the remaining bounded values of
$X$ are absorbed in the implied constant.

The principal character at the chosen modulus is exactly the
character induced by the trivial primitive character, giving the
stated $r_{\chi,\rho}$.
Equivalently one can work at the primitive conductor and remove
the extra Euler factors: the corresponding prime-power towers
contribute only
$O_\chi(\log(q+2)\log(Y+2))$, which is absorbed in the error.

For \eqref{eq:rational-sw-estimate}, take $K=\Q$ and use the ray
character attached to $\rho$, allowing the real infinite place
for odd characters.  Its finite modulus norm is $q$.
The same proof uses $A=B+1$ for $Y\geq\sqrt X$ and the elementary
bound for smaller $Y$.  Finally,
$e^{-c\sqrt{\log X}}\ll_{c,D}(\log X)^{-D}$ for every fixed $D>0$.
\end{proof}


The preceding lemma gives the following consequence, including $Y=0$ by the empty-sum convention.
  For every $B,D>0$, every fixed finite-order
ray character $\chi$ of $E$, and every Dirichlet character $\rho$ of
modulus $q\leq(\log X)^B$,
\begin{equation}\label{eq:gaussian-hecke-sw-input}
 \sum_{\mathrm N\mathfrak a\leq Y}
       \Lambda_E(\mathfrak a)\chi(\mathfrak a)
                         \rho(\mathrm N\mathfrak a)
  = r_{\chi,\rho}Y+O_{H,B,D}(X(\log X)^{-D})
 \quad(0\leq Y\leq X),
\end{equation}
where $r_{\chi,\rho}=1$ if the induced Hecke character
$\chi(\rho\circ\mathrm N)$ is trivial and is zero otherwise.
The omitted Euler factors only change the sum by a polylogarithmic
quantity.  The constant may be ineffective.  We require uniformity in
$\rho,q,Y$, but not in $H$.

\begin{lemma}[Progression density from the fixed-field input]
\label{lem:gaussian-ray-progression}
For $(a,q)=1$ define
$J_q=K\cap\Q(\zeta_q)$ and let $\sigma_a(\zeta_q)=\zeta_q^a$.
For fixed $B,D>0$, uniformly for $q\leq(\log X)^B$ and $0\leq Y\leq X$,
\[
 \sum_{\substack{n\leq Y\\ n\equiv a\pmod q}}\Lambda'_C(n)
 =\frac{\#\{g\in C:g|_{J_q}=\sigma_a|_{J_q}\}}
             {[K\Q(\zeta_q):\Q]}Y
       +O_{H,B,D}(X(\log X)^{-D}).
\]
\end{lemma}
\begin{proof}
Orthogonality of Dirichlet characters and
\eqref{eq:gaussian-ray-fourier} show that the coefficient of $Y$ for
$T_C$ is
\[
 \frac1{\varphi(q)}\sum_{\chi\in\widehat A}c_C(\chi)
       \sum_{\rho\bmod q}\overline{\rho(a)}r_{\chi,\rho}.
\]
The finite sum over $\chi$ has size depending only on $H$, while the
average over $\rho$ cancels its own factor $\varphi(q)$ in the error.

To evaluate this coefficient without a convention about inverse Artin
symbols, put $L=K\Q(\zeta_q)$ and $B_q=\operatorname{Gal}(L/E)$.
Restriction identifies $B_q$ with
\[
 \{(g,s)\in A\times\operatorname{Gal}(\Q(\zeta_q)/\Q):
                       g|_{J_q}=s|_{J_q}\}.
\]
The characters $\chi(g)\rho(s)$ restrict to characters of $B_q$.
Their average on $B_q$ is one precisely when their pullback by Artin
reciprocity is the trivial Hecke character; otherwise it is zero.
This is exactly $r_{\chi,\rho}$.  Expanding $1_C(g)$ and
$1_{s=\sigma_a}$ by Fourier inversion therefore identifies the displayed
coefficient with the proportion of elements of $B_q$ satisfying
$g\in C$ and $s=\sigma_a$.  Its numerator is the numerator in the
lemma and its denominator is $|B_q|=[L:E]$.
Dividing by two using \eqref{eq:gaussian-ray-prime-factor} changes the
denominator to $[L:\Q]$.  The prime-power error is smaller than the stated
error.  This proves the formula.
\end{proof}

Choose a fixed integer $W$ divisible by the conductor of $K^{\mathrm{ab}}$.
It is in particular divisible by $4$.
For a unit $b\pmod W$ put
\[
\begin{aligned}
 \mu_C(b)&=
 \frac{\#\{g\in C:g|_{K^{\mathrm{ab}}}=\sigma_b|_{K^{\mathrm{ab}}}\}}
      {[K:K^{\mathrm{ab}}]},\\
 f_{C,b,W}(n)&=\frac{\varphi(W)}W\Lambda'_C(Wn+b),\\
 \lambda_{b,W}(n)&=\frac{\varphi(W)}W\Lambda'(Wn+b).
\end{aligned}
\]
The prime-supported ordinary von Mangoldt function is denoted by
$\Lambda'$.  Replacing it by the usual $\Lambda$ costs only the same
negligible prime-power error.
The centered weight is
\begin{equation}\label{eq:gaussian-centered-weight}
       h_{C,b,W}=f_{C,b,W}-\mu_C(b)\lambda_{b,W}.
\end{equation}
Here $0\leq\mu_C(b)\leq1$, so
$|h_{C,b,W}(n)|\leq\lambda_{b,W}(n)=O_W(\log(2N))$ on $[N]$.

\begin{corollary}[Centered sums on progressions]
\label{cor:gaussian-centered-progressions}
For every $D>0$,
\[
       \sum_{n\in P}h_{C,b,W}(n)=O_{H,W,D}(N(\log N)^{-D})
\]
uniformly over arithmetic progressions $P$ contained in $[N]$.
\end{corollary}
\begin{proof}
Write $P$ as a residue class modulo $q$ intersected with an interval.
If $q>(\log N)^{D+3}$, its cardinality is at most
$N/(\log N)^{D+3}+1$, and the pointwise logarithmic bound proves the
claim directly.  Otherwise the condition $n\in P$ becomes one residue
class $a\pmod{Wq}$ and an interval for $Wn+b$.

If $(a,Wq)>1$, every prime in this residue class divides $Wq$.  There
are at most $O(\log(Wq))$ such primes, with total logarithmic weight
polylogarithmic in $N$, which is negligible here.  Suppose instead that
$(a,Wq)=1$.
Since $\Q(\zeta_{Wq})$ is abelian and already contains
$K^{\mathrm{ab}}$, one has
$K\cap\Q(\zeta_{Wq})=K^{\mathrm{ab}}$.
Also $a\equiv b\pmod W$, so the lemma's main coefficient is
\[
   \frac{\mu_C(b)}{\varphi(Wq)}.
\]
For the ordinary prime weight, the Siegel--Walfisz main coefficient
is $1/\varphi(Wq)$.  The two terms of
\eqref{eq:gaussian-centered-weight} therefore have equal main terms.
Apply the progression lemma at both interval endpoints, with an error
exponent larger than $D$ and with $X=O_W(N)$, and subtract.  This proves
the claimed bound.  All moduli used are at most a fixed power of
$\log N$, and $W$ is fixed throughout this argument.
\end{proof}


% Standalone section fragment; needs amsmath, amssymb, amsthm.
% Bibliography keys: GTOrbits2012 and GTMobius2012.
\subsection{Transfer from Gaussian ideal norms to polynomial nilsequences}
\label{sec:gn-transfer}

Throughout this subsection $E=\mathbb Q(i)$, $\mathcal O_E=\mathbb Z[i]$,
and $H_0$ is a fixed positive integer. A ray character $\chi$ modulo
$H_0\mathcal O_E$ is extended by zero to ideals not coprime to
$H_0\mathcal O_E$. The trivial character on all ideals is included by
taking $H_0=1$. We use the conventions of \cite{GTOrbits2012}: $G$ is
connected, simply connected and nilpotent, $\Gamma$ is a discrete cocompact
subgroup, and the filtration consists of closed connected rational
subgroups. Lipschitz norms include the supremum norm.

\begin{lemma}[Gaussian norm transfer]\label{lem:gn-norm-transfer}
Fix positive integers $m,d$. There are $\kappa,\sigma>0$, depending only
on $m,d$, and constants $C,X_0,\delta_0>0$, depending on $H_0,m,d$, with
the following property. Suppose
\[
 X\geq X_0,\qquad X^{-\sigma}\leq\delta\leq\delta_0,
 \qquad 1\leq Q\leq\delta^{-\kappa},\qquad
 1\leq q\leq\delta^{-\kappa},\quad q\in\mathbb Z.
\]
Let $G/\Gamma$ have dimension $m$, a filtration $G_\bullet$ of degree at
most $d$, and a $Q$-rational adapted Mal'cev basis. Suppose
$g\in\operatorname{poly}(\mathbb Z,G_\bullet)$ is totally
$\delta$-equidistributed on $[X]$. If $F:G/\Gamma\to\mathbb C$ is
Lipschitz and $\int F=0$, then for every $0\leq A<B\leq X$ and every
$b\in\mathbb Z$,
\begin{equation}\label{eq:gn-transfer}
 \left|\sum_{\substack{A<\mathrm N\mathfrak a\leq B\\
                 \mathrm N\mathfrak a\equiv b\pmod q}}
       \chi(\mathfrak a)F(g(\mathrm N\mathfrak a)\Gamma)\right|
 \leq C X\delta^\kappa\|F\|_{\mathrm{Lip}}.
\end{equation}
There is no lower bound on $B-A$ or on the number of progression terms.
Inserting $\log\mathrm N\mathfrak a$ on the left costs at most an
additional factor $O(\log(2X))$.
\end{lemma}

We prove the lemma using only the one-parameter quantitative Leibman
theorem \cite[Theorem 2.9]{GTOrbits2012}, together with elementary
polynomial and lattice arguments. No multiparameter equidistribution
theorem is used.

\begin{lemma}[Quadratic coefficient extrapolation]\label{lem:gn-quadratic-coefficients}
Let $X\geq1$, $P(t)=a_0+a_1t+a_2t^2\in\mathbb Z[t]$, $a_2\ne0$, and suppose
\[
 |a_i|\leq B_0X^{1-i/2},\qquad
 L\geq X^{1/2}/B_0,\qquad B_0\geq2.
\]
For every scalar polynomial $p:\mathbb Z\to\mathbb R/\mathbb Z$ of
degree at most $d$, there is an integer $1\leq k\leq C_dB_0^{C_d}$ such
that
\[
 \|kp\|_{C^\infty[X]}
   \leq C_dB_0^{C_d}\|p\circ P\|_{C^\infty[L]}.
\]
Here, for $p(n)=\sum_j\alpha_j\binom nj$, we use
$\|p\|_{C^\infty[X]}=\max_{1\leq j\leq d}X^j
\|\alpha_j\|_{\mathbb R/\mathbb Z}$.
\end{lemma}
\begin{proof}
Write $\binom{P(t)}j=\sum_\ell s_{j,\ell}\binom t\ell$. These
coefficients are integers since the left side is integer-valued. Its
degree is $2j$, and comparison of leading coefficients gives
\[
 s_{j,2j}=\frac{(2j)!}{j!}a_2^j\ne0,
 \qquad |s_{j,\ell}|\ll_d B_0^jX^{j-\ell/2}.
\]
For the bound, expand powers of $P$ and change between the monomial and
binomial bases. The change-of-basis matrices have fixed rational
coefficients. Descending induction in $\ell$ shows that every term
contributing to $s_{j,\ell}$ has weighted input degree at least $\ell$,
which gives the displayed estimate.

If $p=\sum_j\alpha_j\binom nj$ and
$p\circ P=\sum_\ell\beta_\ell\binom t\ell$, then
\[
 \beta_{2j}=\sum_{h=j}^d s_{h,2j}\alpha_h.
\]
Put $k_j=\prod_{h=j}^d s_{h,2h}$. Descending induction in $j$ gives
\[
 X^j\|k_j\alpha_j\|_{\mathbb R/\mathbb Z}
   \ll_d B_0^{O_d(1)}\|p\circ P\|_{C^\infty[L]}.
\]
Indeed the $\beta_{2j}$ term is bounded using
$X^jL^{-2j}\leq B_0^{2j}$, and a higher coefficient is multiplied by
$s_{h,2j}=O_d(B_0^hX^{h-j})$, canceling its $X^{-h}$ error.
All multipliers used to clear higher rows are integers of size
$B_0^{O_d(1)}$. Multiplying to the common nonzero integer $|k_1|$ proves
the assertion. The constant coefficient $\alpha_0$ is not used.
\end{proof}

\begin{lemma}[Small horizontal character obstruction]\label{lem:gn-horizontal-obstruction}
There are $c_0,\sigma_0>0$, depending only on $m,d$, such that the
following holds for sufficiently large $X$ and sufficiently small
$\delta\geq X^{-\sigma_0}$. Suppose the basis is $T$-rational and
a nontrivial horizontal character $\eta$ satisfies
\[
 |\eta|\leq T,\qquad \|\eta\circ g\|_{C^\infty[X]}\leq T,
 \qquad T\leq\delta^{-c_0}.
\]
Then $g$ is not totally $\delta$-equidistributed on $[X]$.
\end{lemma}
\begin{proof}
Choose representatives of the nonconstant binomial coefficients of
$\eta\circ g$ satisfying $|\alpha_j|\leq TX^{-j}$. On the initial
interval of length $\lfloor X/(100dT)\rfloor$ its values lie in an arc
of length less than $1/10$, translated by its constant coefficient.
Choose a Lipschitz function on $\mathbb R/\mathbb Z$ which equals one on
that arc, has mean zero and bounded Lipschitz norm. Its composition with
$\eta$ has Lipschitz norm at most $T^{C(m,d)}$, by the Mal'cev metric
estimates in \cite[Appendix A]{GTOrbits2012}. Its Haar mean is zero:
a nontrivial horizontal character maps Haar measure onto Haar measure
on the circle. Choose $c_0$ with $C(m,d)c_0<1/4$ and $c_0<1/4$.
The initial interval is longer than $\delta X$, while the discrepancy
one exceeds $\delta T^{C(m,d)}$. Taking $\sigma_0$ sufficiently small
absorbs the floors. This contradicts total equidistribution.
\end{proof}

\begin{proof}[Proof of Lemma~\ref{lem:gn-norm-transfer}]
Fix a constant $C_*$ which dominates the exponents in
Lemma~\ref{lem:gn-quadratic-coefficients}, the degree-$2d$ case of
\cite[Theorem 2.9]{GTOrbits2012}, and the preceding obstruction.
Choose
\[
 0<\kappa<\min\{(1000C_*^2)^{-1},c_0/(12C_*)\},\qquad
 0<\sigma<\min\{(1000C_*^2)^{-1},\sigma_0\},
 \qquad 2\kappa\sigma<1/4.
\]
These choices depend only on $m,d$. Decrease $\delta_0$, depending on
$H_0,m,d$, whenever a fixed multiplicative constant must be absorbed.

We first record a pullback consequence. Suppose $P$ and $L$ satisfy
Lemma~\ref{lem:gn-quadratic-coefficients} with
$B_0\leq\delta^{-5\kappa}$. Then
$(g(P(t))\Gamma)_{t\in[L]}$ is $\delta^\kappa$-equidistributed.
Otherwise the quantitative Leibman theorem gives a nontrivial $\eta$
with
\[
 |\eta|,\ \|\eta\circ g\circ P\|_{C^\infty[L]}
     \ll\delta^{-C_*\kappa}.
\]
The basis is $\delta^{-\kappa}$-rational as required. The composed
sequence has degree at most $2d$ for the stretched filtration
$\widetilde G_j=G_{\lceil j/2\rceil}$, with the same adapted basis.
This follows from the elementary one-parameter coordinate
characterization in \cite[Lemma 6.7]{GTOrbits2012}.
Extrapolation gives a nonzero integer multiple $k\eta$ with
\[
 |k\eta|,\ \|k\eta\circ g\|_{C^\infty[X]}
       \ll\delta^{-6C_*\kappa},
\]
with $C_*$ chosen sufficiently large initially. It remains nontrivial. The choice
$6C_*\kappa<c_0/2$ and a reduction of $\delta_0$ make this contradict
Lemma~\ref{lem:gn-horizontal-obstruction}. Only ordinary
equidistribution of this pullback is needed.

Every nonzero Gaussian ideal has exactly four generators. The function
$z\mapsto\chi((z))$, extended by zero when $z$ is not coprime to $H_0$,
is periodic modulo $H_0\mathbb Z[i]$ and invariant under units. Thus the
sum in \eqref{eq:gn-transfer} is exactly
\[
 \frac14\sum_{\substack{z=x+iy\ne0\\A<x^2+y^2\leq B\\
                         x^2+y^2\equiv b\pmod q}}
       \chi((z))F(g(x^2+y^2)\Gamma).
\]
The factor $1/4$ also applies to ideals with generators on an axis;
we have introduced no angular fundamental domain.

Set $\rho=\delta^\kappa$ and
$M=\lfloor\rho\sqrt X/H_0\rfloor$. On each horizontal row $y\in
\mathbb Z$, and in each residue class of $x$ modulo $H_0$, partition
the points into consecutive blocks of $M$ points with step $H_0$.
The character is constant on a block. The intersection of the annulus
with a horizontal row consists of at most two intervals, so at most
four blocks per residue class meet an endpoint. There are $O(\sqrt X)$
relevant rows. Consequently the boundary blocks contain
\[
 O(H_0\sqrt X\,M)=O(\rho X)
\]
points. Estimate their contribution trivially. Excluding $z=0$, if
necessary, costs one more bounded summand, absorbed by $\rho X$.

Every remaining contributing block lies wholly in the annulus. Split
its index modulo $q$. The norm congruence is constant on each resulting
class. Such a class has the form
\[
 x=x_0+H_0qn,\quad 1\leq n\leq L,\qquad
 P(n)=(x_0+H_0qn)^2+y^2,
 \qquad L\geq\frac{\rho\sqrt X}{4H_0q}.
\]
The hypotheses on $\delta$ and $\sigma$ ensure $M\geq2q$ and absorb
the floors for large $X$. A harmless endpoint adjustment gives
$|x_0|\leq2\sqrt X$. The coefficients of $P$ satisfy
\[
 |a_0|\ll X,\qquad |a_1|\ll H_0q\sqrt X,\qquad a_2=H_0^2q^2.
\]
Since $H_0$ is fixed and $q\leq\delta^{-\kappa}$, these bounds and the
lower bound for $L$ fit the pullback consequence with
$B_0\leq\delta^{-5\kappa}$ after decreasing $\delta_0$.
Hence the sum on the class is
$O(L\delta^\kappa\|F\|_{\mathrm{Lip}})$. The classes partition the
interior points, of which there are $O(X)$; this elementary bound follows
from the enclosing square. Summing and adding the boundary contribution
proves \eqref{eq:gn-transfer}.

The number of row endpoints is independent of $B-A$, explaining why no
lower interval length is required. Finally apply Abel summation to the
uniform bound with variable upper endpoint to insert $\log\mathrm N
\mathfrak a$ at cost $O(\log(2X))$.
\end{proof}

\begin{remark}\label{rem:gn-product-and-progression}
For Type II, apply the lemma in dimension $2m$ to
$(g_1(n),g_2(n))$ and $F_1\otimes\overline{F_2}$, when this entire
product orbit is totally equidistributed. Mean-zero factors ensure that
the product test function has mean zero; they do not imply
equidistribution of the product orbit. Use the smaller exponents for
dimensions $m$ and $2m$. In dimension zero mean-zero test functions
vanish, so the assertion is immediate.

The constraints $rn\equiv b\pmod q$ and $r'n\equiv b\pmod q$ either
are inconsistent or specify one residue class modulo a divisor of $q$.
Intersecting their interval restrictions only changes $A,B$. Thus all
Type I and Type II truncations are covered by the lemma. If
$\delta=(\log X)^{-B_1}$ and $Q,q\leq(\log X)^{C_1}$, any desired
saving $X(\log X)^{-A_1}$ follows once
$B_1\kappa>\max\{C_1,A_1\}$, with $H_0$ fixed.
\end{remark}


% Depends on the separately proved Lemma lem:gn-norm-transfer.
% No Frei--Tinkova theorem is assumed in this fragment.
% Their ideal-norm Type I/II strategy is retained as methodological attribution.
\subsection{Gaussian ideal Vaughan decomposition and Type I/II estimates}
\label{gi-section}

We give the ideal decomposition and the subsequent integer-parameter
recombination in full.  The architecture is that of
\cite[Section~3]{moebius}, with the Gaussian ideal-norm adaptation
attributed to Frei--Tinkova.  The analytic inputs used here are the
separately proved Lemma~\ref{lem:gn-norm-transfer} and the published
one-parameter results \cite[Theorem~2.9 and Lemma~4.5]{GTOrbits2012}.
In particular, no theorem of Frei--Tinkova is an assumed input.

Fix the Gaussian ray modulus $H_0$.  All ideals in this subsection are
nonzero integral ideals of $\mathbb Z[i]$.  Write $\mathrm N\mathfrak a$
for the absolute ideal norm, $\mu_E$ for ideal Mobius, and
$\Lambda_E(\mathfrak p^j)=\log\mathrm N\mathfrak p$ for $j\geq1$, with
$\Lambda_E$ zero on other ideals.  A ray character $\chi$ is extended by
zero on ideals not coprime to $H_0$.  It is completely multiplicative
and satisfies $|\chi|\leq1$.  We use a Lipschitz norm that includes the
supremum norm.  All nilmanifolds have connected, simply connected
nilpotent covering group, a cocompact lattice, and an adapted rational
Mal'cev basis, with the filtration and polynomial-sequence conventions
of \cite{GTOrbits2012}.

\begin{lemma}[Ideal multiplicities and moments]\label{gi-moments}
Let $a_E(n)=\#\{\mathfrak a:\mathrm N\mathfrak a=n\}$.  For $X\geq1$,
\begin{align}
 a_E(n)&\leq\tau(n),\label{gi-multiplicity}\\
 \sum_{n\leq X}a_E(n)^2&\leq X(1+\log X)^3,&
 \sum_{n\leq X}\frac{a_E(n)^2}{n}&\leq(1+\log X)^4.\label{gi-two-moments}
\end{align}
Moreover
\begin{equation}\label{gi-ideal-count}
 A_E(X):=\#\{\mathfrak a:\mathrm N\mathfrak a\leq X\}
       =\frac{\pi}{4}X+O(\sqrt X+1).
\end{equation}
In particular $A_E(X)\ll X$, and the number of ideals of norm in
$(B/2,B]$ is comparable to $B$ for all sufficiently large $B$.
\end{lemma}
\begin{proof}
At a split rational prime $p$ one has $a_E(p^k)=k+1$; at an inert prime
it is one for even $k$ and zero for odd $k$; and $a_E(2^k)=1$.
Multiplicativity proves \eqref{gi-multiplicity}.  Prime by prime,
$(k+1)^2\leq\binom{k+3}{3}$, so $\tau(n)^2\leq d_4(n)$.
Summing over ordered quadruples with product at most $X$, and summing
the last coordinate first, bounds $\sum_{n\leq X}d_4(n)$ by
$X(1+\log X)^3$.  Bounding each of the four harmonic sums separately
gives $\sum_{n\leq X}d_4(n)/n\leq(1+\log X)^4$.
Every ideal has exactly four Gaussian generators.  Counting lattice
points in the disk of radius $\sqrt X$, and dividing by four, proves
\eqref{gi-ideal-count}: the unit squares meeting its boundary lie
between circles of radii $\sqrt X-\sqrt2$ and $\sqrt X+\sqrt2$ and
number $O(\sqrt X+1)$.
\end{proof}

\begin{lemma}[Exact ideal Vaughan identity]\label{gi-vaughan}
Let $N\geq8$, $U=V=N^{1/3}$, $L=\log(2N)$, and let
$f:\mathbb N\longrightarrow\mathbb C$ satisfy $|f|\leq1$ and vanish
outside $(N/2,N]$.  Then
\begin{equation}\label{gi-vaughan-sum}
 \sum_{\mathfrak a}\Lambda_E(\mathfrak a)\chi(\mathfrak a)
                       f(\mathrm N\mathfrak a)=T_0-T_1+T_2,
\end{equation}
where
\begin{align*}
 T_0&=\sum_{\mathrm N\mathfrak d\leq U}\mu_E(\mathfrak d)\chi(\mathfrak d)
       \sum_{\mathfrak b}\log(\mathrm N\mathfrak b)\chi(\mathfrak b)
                  f(\mathrm N\mathfrak d\,\mathrm N\mathfrak b),\\
 T_1&=\sum_{\mathrm N\mathfrak d\leq UV}\alpha_{\mathfrak d}\chi(\mathfrak d)
       \sum_{\mathfrak b}\chi(\mathfrak b)
                  f(\mathrm N\mathfrak d\,\mathrm N\mathfrak b),\\
 T_2&=\sum_{V<\mathrm N\mathfrak d\leq N/U}\beta_{\mathfrak d}\chi(\mathfrak d)
       \sum_{\mathrm N\mathfrak b>U}\mu_E(\mathfrak b)\chi(\mathfrak b)
                  f(\mathrm N\mathfrak d\,\mathrm N\mathfrak b),
\end{align*}
with
\begin{align*}
 \alpha_{\mathfrak d}
   &=\sum_{\substack{\mathfrak b\mathfrak c=\mathfrak d\\
              \mathrm N\mathfrak b\leq U,\ \mathrm N\mathfrak c\leq V}}
                        \mu_E(\mathfrak b)\Lambda_E(\mathfrak c),&
 \beta_{\mathfrak d}
   &=\sum_{\substack{\mathfrak c\mid\mathfrak d\\\mathrm N\mathfrak c>V}}
                        \Lambda_E(\mathfrak c).
\end{align*}
Both coefficients have absolute value at most $L$.
\end{lemma}
\begin{proof}
Convolution on ideals satisfies
$\mu_E*1=\mathbf1_{\mathfrak a=(1)}$, $1*\Lambda_E=\log\mathrm N$,
and $\mu_E*\log\mathrm N=\Lambda_E$, as one checks on prime powers.
Splitting $\mu_E$ at $U$ and $\Lambda_E$ at $V$ gives the exact identity
\begin{equation}\label{gi-convolution-identity}
 \Lambda_E=\mu_{\leq U}*\log\mathrm N+\Lambda_{\leq V}
       -\mu_{\leq U}*1*\Lambda_{\leq V}
       +\mu_{>U}*1*\Lambda_{>V}.
\end{equation}
The second term vanishes on the support of $f$.  Complete multiplicativity
of $\chi$, including its zero extension, permits separation of each
convolution factor and yields \eqref{gi-vaughan-sum} with exactly the
stated cutoffs.  Finally
$\sum_{\mathfrak c\mid\mathfrak d}\Lambda_E(\mathfrak c)
=\log\mathrm N\mathfrak d$ and every summand is nonnegative.  This
proves the coefficient bounds, with no discarded endpoint or error term.
\end{proof}

\begin{proposition}[Quantitative extraction]\label{gi-extraction}
There are absolute constants $c_0>0$ and $N_0$ such that the following
holds.  Under Lemma~\ref{gi-vaughan}, suppose $N\geq N_0$,
$0<\epsilon\leq1$, and
\begin{equation}\label{gi-large-correlation}
 \left|\sum_{\mathfrak a}\Lambda_E(\mathfrak a)\chi(\mathfrak a)
                         f(\mathrm N\mathfrak a)\right|\geq\epsilon NL.
\end{equation}
Put $\theta=c_0\epsilon^4/L^{24}$.  At least one of the following holds.
\begin{enumerate}
\item Some dyadic $R\leq2N^{1/3}$ has at least $\theta R$ integers
$r\in(R/2,R]$ satisfying
\[
 \left|\sum_{\mathfrak b}\log(\mathrm N\mathfrak b)\chi(\mathfrak b)
                   f(r\mathrm N\mathfrak b)\right|\geq\theta N/r.
\]
\item Some dyadic $R\leq2N^{2/3}$ has at least $\theta R$ integers
$r\in(R/2,R]$ satisfying the same inequality without the logarithm.
\item There are dyadic $R,B$ with
\begin{equation}\label{gi-dyadic-ranges}
 N^{1/3}<R\leq2N^{2/3},\qquad N/2<RB<4N,
\end{equation}
and at least $\theta R^2$ pairs $r,r'\in(R/2,R]$ satisfying
\begin{equation}\label{gi-type-two-correlation}
 \left|\sum_{B/2<\mathrm N\mathfrak b\leq B}
       f(r\mathrm N\mathfrak b)\overline{f(r'\mathrm N\mathfrak b)}
       \right|\geq\theta B.
\end{equation}
\end{enumerate}
\end{proposition}
\begin{proof}
One of $|T_i|$ is at least $\epsilon NL/3$.  For $T_0$ put
\[
 G(r)=\frac{r}{NL}\left|\sum_{\mathfrak b}\log(\mathrm N\mathfrak b)
                   \chi(\mathfrak b)f(r\mathrm N\mathfrak b)\right|;
\]
for $T_1$ use $r/N$ times the absolute unweighted inner sum.  In either
case Lemma~\ref{gi-moments} gives $G\ll1$ and
\[
 \epsilon\ll\sum_{r\leq N^e}\frac{a_E(r)G(r)}r,
 \qquad e=1/3\text{ or }2/3.
\]
Cauchy--Schwarz and \eqref{gi-two-moments} imply
$\sum_{r\leq N^e}G(r)^2/r\gg\epsilon^2/L^4$.
There are $O(L)$ dyadic cells, so one satisfies
\[
 \sum_{R/2<r\leq R}G(r)^2\gg R\epsilon^2/L^5.
\]
Since $G\ll1$, at least $cR\epsilon^2/L^5$ terms have
$G(r)\geq c\epsilon/L^{5/2}$.  Both the count and the claimed amplitudes
follow after decreasing $c_0$.

For $T_2$, split both ideal norms into dyadic cells.  There are $O(L^2)$
rectangles, and a nonzero cell satisfies \eqref{gi-dyadic-ranges} and
$B\geq N^{1/3}/4$.  A rectangle contributes at least $c\epsilon N/L$;
we use the weaker bound $c\epsilon N/L^2$.  Grouping $\mathfrak d$ by
$r=\mathrm N\mathfrak d$ writes this contribution as
\[
 T=\sum_{R/2<r\leq R}A_r\sum_{\mathfrak b\in\mathcal B}c_{\mathfrak b}
                   f(r\mathrm N\mathfrak b),\qquad
 \mathcal B=\{\mathfrak b:B/2<\mathrm N\mathfrak b\leq B\}.
\]
Here $|A_r|\leq L a_E(r)$ and $|c_{\mathfrak b}|\leq1$; the restriction
$\mathrm N\mathfrak b>U$ belongs to $c_{\mathfrak b}$.  All product
cutoffs remain in $f$, so the domain is rectangular.  By
\eqref{gi-two-moments}, $\sum_r|A_r|^2\ll RL^5$.  Thus, with
$F_{r,\mathfrak b}=f(r\mathrm N\mathfrak b)$,
\begin{equation}\label{gi-first-cauchy}
 T_*:=\sum_r\left|\sum_{\mathfrak b\in\mathcal B}
                          c_{\mathfrak b}F_{r,\mathfrak b}\right|^2
          \gg \epsilon^2 RB^2/L^9.
\end{equation}
Put $M=|\mathcal B|\asymp B$.  Expansion and Cauchy--Schwarz in the
ordered pair $(\mathfrak b,\mathfrak b')$ give
\begin{align}
 T_*^2
 &=\left|\sum_{\mathfrak b,\mathfrak b'}
 c_{\mathfrak b}\overline{c_{\mathfrak b'}}
 \sum_r F_{r,\mathfrak b}\overline{F_{r,\mathfrak b'}}\right|^2\notag\\
 &\leq M^2\sum_{\mathfrak b,\mathfrak b'}
       \left|\sum_rF_{r,\mathfrak b}\overline{F_{r,\mathfrak b'}}\right|^2
 =M^2\sum_{r,r'}\left|\sum_{\mathfrak b}
                  F_{r,\mathfrak b}\overline{F_{r',\mathfrak b}}\right|^2.
 \label{gi-second-cauchy}
\end{align}
Consequently
\[
 \mathbb E_{r,r'}\left|\mathbb E_{\mathfrak b\in\mathcal B}
              F_{r,\mathfrak b}\overline{F_{r',\mathfrak b}}\right|^2
       \gg\epsilon^4/L^{18}.
\]
Every summand is at most one.  Thresholding gives
\eqref{gi-type-two-correlation} and its stated count, with the weaker
common parameter $\theta=c_0\epsilon^4/L^{24}$.  The character and Mobius
coefficients have disappeared.  All conjugates in
\eqref{gi-second-cauchy} are necessary.
\end{proof}

\paragraph{Progressions and interval restrictions.}\label{gi-progressions}
In the application,
\[
 f(n)=1_{N/2<n\leq N}\,1_{u<n\leq v}\,1_{n\equiv b\pmod q}
                              F(g(n)\Gamma).
\]
For Type I the congruence $rn\equiv b\pmod q$ is either insoluble or
one class modulo $q/\gcd(r,q)$.  In Type II its intersection with
$r'n\equiv b\pmod q$ is either empty or one class modulo
\[
 \operatorname{lcm}\bigl(q/\gcd(r,q),q/\gcd(r',q)\bigr),
\]
which divides $q$.  The remaining restrictions are interval intersections.
Lemma~\ref{lem:gn-norm-transfer} covers them without a lower bound on
their lengths.  On the product nilmanifold the function is
$F(x)\overline{F(y)}$, whose Haar integral is
$|\int F|^2=0$.  After normalizing $\|F\|_{\mathrm{Lip}}=1$, its
Lipschitz norm for the sum metric is at most two.  Thus all mean-zero,
modulus and interval hypotheses for product norm transfer are satisfied.

\begin{lemma}[Integer-parameter recombination]\label{gi-recombination}
For each $m,d$ there are $C_{\rm rec},K_{\rm rec}\geq2$ with the following
property.  Let $0<\eta<\eta_0$, $N\geq\eta^{-K_{\rm rec}}$, and let $g$
have degree at most $d$ on an $m$-dimensional nilmanifold with an
$\eta^{-1}$-rational adapted basis.  Each of the following implies that
$(g(n)\Gamma)_{n\leq N}$ is not totally $\eta^{C_{\rm rec}}$-equidistributed.
\begin{enumerate}
\item For at least $\eta R$ values $r\in(R/2,R]$, where
$R\leq2N^{2/3}$, the orbit $(g(rn)\Gamma)_{n\leq\lceil N/r\rceil}$
is not totally $\eta$-equidistributed.
\item For at least $\eta R^2$ pairs $r,r'\in(R/2,R]$, where
$N^{1/3}/4\leq R\leq2N^{2/3}$ and $N/2<RB<4N$, the product orbit
$(g(rn),g(r'n))_{n\leq B}$ is not totally $\eta$-equidistributed.
\end{enumerate}
\end{lemma}
\begin{proof}
All constants in this proof depend only on $m,d$.  Failure of total
equidistribution gives a progression $a+q[\ell]\subseteq[M]$ with
$\ell\geq\eta M$, $q\leq2/\eta$, and $|a|\leq2M$ on which ordinary
$\eta$-equidistribution fails.  Apply
\cite[Theorem~2.9]{GTOrbits2012}, in dimension $m$ or $2m$.  A nonzero
horizontal character has magnitude and restricted smoothness norm at
most $\eta^{-C_1}$, after shrinking $\eta_0$ to absorb fixed factors.

Here is the required removal of a possibly large shift.  Write the scalar
phase $p(a+qt)$ in monomials.  Multiplication by $d!$ clears conversion
from its binomial coefficients.  Inverting $t=(n-a)/q$ and multiplying
by $q^d$ clears every remaining denominator.  The error in the coefficient
of $n^i$ is bounded by
\begin{equation}\label{gi-shift-removal}
 C_d\eta^{-C_1}\sum_{j=i}^d q^{d-j}|a|^{j-i}\ell^{-j}
       \leq\eta^{-C_2}M^{-i},\qquad i\geq1.
\end{equation}
The integer parts disappear modulo one.  Thus a nonzero integer multiple
of the character has magnitude and $C^\infty[M]$ norm at most
$\eta^{-C_2}$.  This argument does not treat $a=O(M)$ as a bounded-height
rational parameter.  There are at most $\eta^{-C_3}$ possible resulting
characters, since their lattice has rank at most $2m$.  Pigeonholing
therefore fixes one character on a set of density at least $\eta^{C_4}$.

In Type I write
$\psi(g(n))=\sum_{j=0}^d\beta_jn^j$ modulo one.  For the selected set
$S\subseteq(R/2,R]$, coefficient conversion gives
\begin{equation}\label{gi-one-coefficients}
 \|d!\beta_jr^j\|_{\mathbb R/\mathbb Z}
       \leq\eta^{-C_5}(R/N)^j,\qquad 1\leq j\leq d.
\end{equation}
If $R\leq\eta^{-C_6}$ for a sufficiently large fixed $C_6$, choose one
$r\in S$ and multiply the character by $d!r^d$.  Multiplication of the
$j$th bound by $r^{d-j}$ gives magnitude and $C^\infty[N]$ norm bounded
by a fixed power of $\eta^{-1}$.  This treats every small dyadic scale,
including $R=1$.

Otherwise apply \cite[Lemma~4.5]{GTOrbits2012} to
$r\mapsto d!\beta_jr^j$ on $[\lceil R\rceil]$.  The selected set has
density at least a fixed power of $\eta$, and its values lie in an arc
of length at most $2\eta^{-C_5}(R/N)^j$.  Taking $K_{\rm rec}$ large,
the bound $R\leq2N^{2/3}$ makes this length at most half the density.
The choice of $C_6$ also places the recurrence length beyond the fixed
inverse density powers.  The published recurrence lemma supplies
nonzero integers $t_j$, $|t_j|\leq\eta^{-C_7}$, such that
\begin{equation}\label{gi-recurrent-coefficients}
 \|t_j(d!)^2\beta_j\|_{\mathbb R/\mathbb Z}
                    \leq\eta^{-C_8}N^{-j}.
\end{equation}
The extra factorial clears the leading-coefficient conversion.  Taking
the product of the $t_j$ and the fixed factorials gives a single nonzero
horizontal character of magnitude and $C^\infty[N]$ norm at most
$\eta^{-C_9}$.

For Type II the pigeonholed character is $\psi_1\oplus\psi_2$, not both
zero.  Write the two scalar phases with coefficients $\beta_j,\beta'_j$.
The information replacing \eqref{gi-one-coefficients} is
\begin{equation}\label{gi-two-coefficients}
 \|d!(\beta_jr^j+\beta'_jr'^j)\|_{\mathbb R/\mathbb Z}
                        \leq\eta^{-C_5}(R/N)^j.
\end{equation}
Interchange the two coordinates if necessary so that $\psi_1\ne0$.
Fix an $r'$ with at least $\eta^{C_4}R$ partners.  Equation
\eqref{gi-two-coefficients} puts $d!\beta_jr^j$ in an arc centered at
$-d!\beta'_jr'^j$.  The same published recurrence lemma permits this
arbitrary center and gives \eqref{gi-recurrent-coefficients} for
$\psi_1$.  Its large-length requirement follows from
$R\geq N^{1/3}/4$ and a sufficiently large $K_{\rm rec}$.  There is no
substitution of $(r-r_0)^j$ for $r^j-r_0^j$.

Finally let $\Psi$ denote the nonzero character just obtained.  Its
binomial expansion shows that for
$1\leq n\leq c\eta^{2C_9+4}N$ the phase $\Psi(g(n))$ stays within
$1/10$ of $\Psi(g(0))$.  Choose a mean-zero Lipschitz circle function
equal to one on that arc and compose it with $\Psi$.  Its mean is zero
because a nonzero horizontal character pushes Haar measure to circle
Haar measure.  Its Lipschitz norm is bounded by a fixed power of
$\eta^{-1}$ by \cite[Appendix~A]{GTOrbits2012}.  Choose $C_{\rm rec}$
larger than the length and Lipschitz exponents, and then enlarge
$K_{\rm rec}$ so the initial interval is nonempty.  Its average one
contradicts total $\eta^{C_{\rm rec}}$-equidistribution.  This proves
both assertions.  All exponent constants are finite combinations of
the published one-parameter constants and degree-$d$ polynomial algebra.
\end{proof}

\begin{proposition}[Power-form Gaussian ideal minor-arc bound]
\label{gi-minor-arc}
For each $m,d$ there is $c_{\rm G}>0$ with the following property.
Fix $H_0$ and a finite-order Gaussian ray character $\chi$ of modulus
$H_0$, extended by zero.  Let $G/\Gamma$ have dimension $m$, filtration
degree $d$, and a $Q$-rational adapted basis, with $Q\geq2$.  Suppose
$(g(n)\Gamma)_{n\leq N}$ is totally $\delta$-equidistributed, where
$0<\delta<1$.  If $F$ is Lipschitz and $\int F=0$, then for every
interval progression $P\subseteq[N]$ of common difference at most $Q$,
\begin{equation}\label{gi-power-bound}
 \left|\sum_{\mathrm N\mathfrak a\leq N}\Lambda_E(\mathfrak a)
       \chi(\mathfrak a)1_P(\mathrm N\mathfrak a)
       F(g(\mathrm N\mathfrak a)\Gamma)\right|
 \ll_{H_0,m,d} QN\log(2N)^7\delta^{c_{\rm G}}\|F\|_{\mathrm{Lip}}.
\end{equation}
No lower bound on the length of $P$ or on $\delta$ is required.
\end{proposition}
\begin{proof}
Normalize $\|F\|_{\mathrm{Lip}}=1$.  Take the smaller positive exponents
$\kappa,\sigma$ in Lemma~\ref{lem:gn-norm-transfer} for dimensions $m$
and $2m$, and decrease $\kappa$ to at most one.  Partial summation gives
the same transfer bound with an extra factor $\log(2M)$ for a
$\log\mathrm N\mathfrak b$ weight: the unweighted lemma holds for
every partial interval, and $\int_1^Mdt/t\leq\log(2M)$.
Let $C_{\rm rec},K_{\rm rec}$ be as in Lemma~\ref{gi-recombination},
and choose
\begin{equation}\label{gi-exponent-choice}
 a_*=\frac1{4C_{\rm rec}},\qquad
 c_* =\frac{a_*\kappa}{32},\qquad
 s_* =\min\left\{\frac1{100},\frac1{4a_*K_{\rm rec}},
                             \frac{\sigma}{12a_*}\right\}.
\end{equation}
First assume $N^{-s_*}\leq\delta\leq\delta_0$, with the fixed cutoff
$\delta_0$ sufficiently small, and restrict the sum to $(N/2,N]$.
We claim that its absolute value is
\begin{equation}\label{gi-shell-bound}
                  \ll QNL^7\delta^{c_*},\qquad L=\log(2N).
\end{equation}
The elementary bound is $O(NL)$.  Thus the claim is immediate when
$QL^6\delta^{c_*}\geq1$, and in the other case
$Q\leq\delta^{-c_*}$ and $L\leq\delta^{-c_*/6}$.
If \eqref{gi-shell-bound} fails with a sufficiently large implied
constant, Proposition~\ref{gi-extraction} applies with
$\epsilon$ a fixed large multiple of $QL^6\delta^{c_*}$; if this
exceeds the permitted range, the elementary bound already proves the
claim.  The extracted density and normalized amplitude are at least
a fixed positive constant times $\delta^{4c_*}$, hence at least
$\delta^{5c_*}$ after decreasing $\delta_0$.

Set $\eta=\delta^{a_*}$.  If an extracted orbit were totally
$\eta$-equidistributed, Lemma~\ref{lem:gn-norm-transfer}, with the
progression analysis above, would bound the extracted sum by
\[
 O(M\delta^{a_*\kappa}L^e),\qquad e=0\text{ or }1.
\]
Here $M=\lceil N/r\rceil\geq N^{1/3}/2$ in Type I and
$M=B\geq N^{1/3}/4$ in Type II.  The rationality and modulus hypotheses
hold because $c_*<a_*\kappa$.  The lower-parameter hypothesis holds
because $a_*s_*\leq\sigma/12$; increasing the fixed threshold for $N$
absorbs the constants two and four.  Since $a_*\kappa=32c_*$, the
displayed bound is smaller than $M\delta^{5c_*}$, even for $e=1$.
This is a contradiction.  Thus every extracted orbit fails total
$\eta$-equidistribution.  Its density is at least $\eta$, since
$a_*>5c_*$.  Also $N\geq\eta^{-K_{\rm rec}}$ by
$a_*s_*K_{\rm rec}\leq1/4$, and $Q\leq\eta^{-1}$.
Lemma~\ref{gi-recombination} now contradicts total
$\delta$-equidistribution, because
$\eta^{C_{\rm rec}}=\delta^{1/4}>\delta$.  This proves
\eqref{gi-shell-bound}.  Fixed values $\delta\geq\delta_0$ are covered
by the elementary bound with a larger constant.

We spell out the passage through all scales.  Use exact real endpoints
$M_j=2^{-j}N$ and retain shells with $M_j\geq2\sqrt\delta N$.
The discarded norms are at most $2\sqrt\delta N$ and their total
absolute weight is $O(\sqrt\delta NL)$ by Lemma~\ref{gi-moments}.
There is no rounding of ideal-norm cutoffs.  On $[\lfloor M_j\rfloor]$
the orbit is totally $2\sqrt\delta$-equidistributed: every progression
of the requisite length has at least $\delta N$ terms for sufficiently
large $N$, and the original discrepancy bound is stronger.
For $\delta\geq N^{-s_*/4}$ the lower-parameter condition for the shell
bound holds on every retained scale, since
$M_j\geq N^{1-s_*/8}$ up to fixed constants.  The progression is merely
intersected with the shell; its common difference remains at most $Q$.
Summing the shell estimates and using $\sum_jM_j\leq2N$ gives
\begin{equation}\label{gi-full-intermediate}
                  O(QNL^7\delta^{c_*/2}).
\end{equation}
For $1/N<\delta<N^{-s_*/4}$ replace $\delta$ by
$\delta^{s_*/8}\geq N^{-s_*/8}$.  Total equidistribution at the
original parameter implies it at this larger one.  The exponent in
\eqref{gi-full-intermediate} becomes $c_*s_*/16$.
For $\delta\leq1/N$, singleton progressions imply
$|F(g(n)\Gamma)|\leq\delta$ at each integer $n$, so the elementary
$O(NL)$ weight bound suffices.  All remaining bounded $N$ are absorbed
by the fixed constant.  Thus one may take
\begin{equation}\label{gi-final-exponent}
                         c_{\rm G}=c_*s_*/16>0.
\end{equation}
Its dependence is only on $m,d$ and the corresponding norm-transfer
exponents; $H_0$ enters constants and thresholds, not this exponent.
\end{proof}

\begin{corollary}[Ordinary prime term]\label{gi-ordinary-primes}
Under the nilmanifold hypotheses of Proposition~\ref{gi-minor-arc},
there is $c_{\mathbb Q}>0$, depending only on $m,d$, such that
\begin{equation}\label{gi-ordinary-bound}
 \left|\sum_{n\leq N}\Lambda'(n)1_P(n)F(g(n)\Gamma)\right|
 \ll_{m,d} QN\log(2N)^7\delta^{c_{\mathbb Q}}\|F\|_{\mathrm{Lip}},
\end{equation}
where $\Lambda'(p)=\log p$ on rational primes and is zero otherwise.
The same bound holds for ordinary von Mangoldt $\Lambda$.
\end{corollary}
\begin{proof}
For positive integers the convolution identities
$\mu*1=1_{n=1}$, $1*\Lambda=\log$, and $\mu*\log=\Lambda$ give
\eqref{gi-convolution-identity} with exactly the displayed norm cutoffs
replaced by integer cutoffs.  Its coefficients are bounded by
$\sum_{c\mid d}\Lambda(c)=\log d$.  Integer multiplicity is one,
so both moment bounds in Lemma~\ref{gi-moments} remain valid.  The
two Cauchy--Schwarz calculations \eqref{gi-first-cauchy} and
\eqref{gi-second-cauchy} apply to the resulting finite integer sums
and give Proposition~\ref{gi-extraction} with the same conservative
parameter $c_0\epsilon^4/L^{24}$.  Lemma~\ref{gi-recombination} concerns
only integer dilation parameters and is unchanged.

The unweighted replacement for norm transfer follows directly from
total equidistribution.  If an interval progression $J\subseteq[X]$
has $|J|\geq\zeta X$, then
\[
 \left|\sum_{n\in J}F(g(n)\Gamma)\right|
       \leq |J|\zeta\|F\|_{\mathrm{Lip}}
       \leq X\zeta\|F\|_{\mathrm{Lip}}.
\]
If $|J|<\zeta X$, its cardinality and the supremum norm give the same
bound.  This treats every modulus and every truncated interval.
Partial summation supplies the logarithmic weight; the identical
long-progression/short-progression dichotomy treats the mean-zero
product function $F\otimes\overline F$ on the product nilmanifold.
Thus the exponent choices and all shell steps in the proof of
Proposition~\ref{gi-minor-arc} apply, with any sufficiently small
fixed positive transfer exponents, to give \eqref{gi-ordinary-bound}
for $\Lambda$.  No extension of a Mobius theorem to $\Lambda$ is assumed.

Finally proper prime powers contribute at most
$O(\sqrt N\log(2N)^2\|F\|_{\mathrm{Lip}})$, by counting at most
$N^{1/k}$ possible bases for each $2\leq k\leq\log_2N$.
Shrink $c_{\mathbb Q}$ below $1/4$.  If $\delta>1/N$, this error is
absorbed because $N\delta^{c_{\mathbb Q}}\geq N^{1-c_{\mathbb Q}}$.
If $\delta\leq1/N$, singleton equidistribution and
$\sum_{n\leq N}\Lambda'(n)\leq N\log(2N)$ prove the bound directly.
This gives the assertion for $\Lambda'$ as well.
\end{proof}



% Section fragment. Requires the gn-* lemmas from norm_transfer.tex.
\subsection{Fixed-modulus affine interpolation}
\label{sec:ga-affine}

Fix integers $W\geq1$ and $0\leq b<W$. The modulus $W$ remains fixed
as the size parameter tends to infinity. We first justify the affine
interpolation needed to apply ideal-norm estimates to a $W$-tricked
sequence.

\begin{lemma}[Linear coefficient extrapolation]\label{lem:ga-linear-coefficients}
Let $X\geq1$ and let $p:\mathbb Z\to\mathbb R/\mathbb Z$ have degree at most $d$.
Suppose $P(t)=u+vt$ with $u,v\in\mathbb Z$, $v\ne0$,
$|u|\leq C X$, $|v|\leq E$, and $L\geq X/E$, where $E\geq2$.
There is an integer $1\leq k\leq C_dE^{C_d}$ such that
\[
 \|kp\|_{C^\infty[X]}
 \leq C_d(C+E)^{C_d}\|p\circ P\|_{C^\infty[L]}.
\]
\end{lemma}
\begin{proof}
Expand $\binom{u+vt}j$ in the binomial basis in $t$. All coefficients
are integers, the diagonal coefficient is $v^j$, and the coefficient in
row $\ell$ is $O_d((C+E)^jX^{j-\ell})$. Starting in degree $d$, clear
the nonzero diagonal entries successively. Their product is a nonzero
integer bounded by $E^{O_d(1)}$. At degree $j$ use
$X^j/L^j\leq E^j$; a higher-row error is multiplied by
$O_d((C+E)^hX^{h-j})$, canceling its $X^{-h}$ factor. Multiplication
to the common product of diagonal entries proves the claim, exactly
as in Lemma~\ref{lem:gn-quadratic-coefficients}.
\end{proof}

\begin{lemma}[Affine interpolation preserves quantitative distribution]
\label{lem:ga-affine-interpolation}
Fix dimension and degree bounds $m,d$. There are $a,\sigma>0$, depending
only on $m,d$, and thresholds depending on $W,m,d$, with the following
property. Suppose a nilmanifold and its rational filtration satisfy the
Green--Tao conventions, its adapted basis is $Q$-rational, and
$g\in\operatorname{poly}(\mathbb Z,G_\bullet)$. Suppose also that
\[
 N^{-\sigma}\leq\delta\leq\delta_0(W,m,d),\qquad
 Q\leq\delta^{-a},
\]
and that $g$ is totally $\delta$-equidistributed on $[N]$. There is
$\widetilde g\in\operatorname{poly}(\mathbb Z,G_\bullet)$ satisfying
\[
 \widetilde g(Wn+b)=g(n)\quad(n\in\mathbb Z)
\]
which is totally $\delta^a$-equidistributed on $[WN+b]$.
\end{lemma}
\begin{proof}
Use the one-parameter coordinate characterization in
\cite[Lemma 6.7]{GTOrbits2012}. Each second-kind Mal'cev coordinate of
$g(n)$ is a polynomial, with the prescribed degree bound from the
filtration. Evaluate those same polynomials at real arguments and apply
the inverse coordinate map, obtaining $g_{\mathbb R}(t)$. Set
\[
 \widetilde g(t)=g_{\mathbb R}((t-b)/W).
\]
Affine composition preserves the coordinate degree bounds, so the same
characterization shows that this is a polynomial sequence for the
original filtration. It gives the stated identity on integer $n$.
Equivalently, one may use the Taylor product with real powers; those
powers exist in a connected simply connected nilpotent group.

Put $X=WN+b$ and $\epsilon=\delta^a$. Suppose the asserted total
$\epsilon$-equidistribution fails. Then on a progression
\[
 u+v[L]\subseteq[X],\qquad L\geq\epsilon X,
 \quad u,v\in\mathbb Z,\quad
 1\leq v\leq2/\epsilon,\quad |u|\leq2X,
\]
the sequence is not $\epsilon$-equidistributed. By the one-parameter
quantitative Leibman theorem \cite[Theorem 2.9]{GTOrbits2012}, there is
a nontrivial horizontal character $\eta$ satisfying
\[
 |\eta|\ll\epsilon^{-C},\qquad
 \|\eta\circ\widetilde g\circ(u+vt)\|_{C^\infty[L]}
       \ll\epsilon^{-C}.
\]
The rationality hypothesis holds since $Q\leq\epsilon^{-1}$.
Lemma~\ref{lem:ga-linear-coefficients} gives a positive integer
$k\ll\epsilon^{-C}$ such that
\[
 \|k\eta\circ\widetilde g\|_{C^\infty[X]}
      \ll\epsilon^{-C}.
\]
Forward composition with the fixed integer polynomial $Wn+b$ gives
\[
 \|k\eta\circ g\|_{C^\infty[N]}
      \ll_{W,d}\epsilon^{-C}.
\]
For this last bound expand $\binom{Wn+b}j$ in the binomial basis; its
matrix is integral, and $X$ is comparable to $WN$. The character
$k\eta$ remains nontrivial.

Let $C_{\mathrm{aff}}$ bound all exponents in these estimates, including
the character size. Fix explicitly
\[
 0<a<\min\{1,c_0/(2C_{\mathrm{aff}})\},\qquad
 0<\sigma<\min\{\sigma_0,1/(2a)\},
\]
where $c_0,\sigma_0$ are from
Lemma~\ref{lem:gn-horizontal-obstruction}. Decrease $\delta_0$ to
absorb fixed $W$-dependent constants. The resulting character bounds
contradict that lemma and total $\delta$-equidistribution of $g$.
The inequality $a\sigma<1/2$ ensures $\epsilon X\gg N^{1/2}$, so all
progression parametrizations and floor estimates used above hold for
large $N$. This proves the assertion on the entire interval $[WN+b]$,
including progressions outside the residue class $b$ modulo $W$.
\end{proof}

The Gaussian input in the following lemma is furnished by
Proposition~\ref{gi-minor-arc}; its rational counterpart is
Corollary~\ref{gi-ordinary-primes}.

\begin{lemma}[$W$-tricked minor-arc interface]\label{lem:ga-wtrick-ma}
For $E=\mathbb Q(i)$ and a fixed ray character $\chi$, put
\[
 T_\chi(t)=\sum_{\mathrm N\mathfrak a=t}
             \Lambda_E(\mathfrak a)\chi(\mathfrak a).
\]
Assume the mean-zero ideal estimate
\begin{equation}\label{eq:ga-ideal-input}
 \left|\sum_{\substack{\mathrm N\mathfrak a\leq X\\
                          \mathrm N\mathfrak a\in R}}
  \Lambda_E(\mathfrak a)\chi(\mathfrak a)
             F_0(h(\mathrm N\mathfrak a)\Gamma)\right|
 \ll Q_1X\log(2X)^7\eta^{c_E}\|F_0\|_{\mathrm{Lip}}
\end{equation}
for a positive $c_E$, whenever $h$ is totally $\eta$-equidistributed,
the basis is $Q_1$-rational, and $R$ is any interval progression of
modulus at most $Q_1$. Then, under the hypotheses of
Lemma~\ref{lem:ga-affine-interpolation}, for every Lipschitz $F_0$
with Haar mean zero and every interval progression
$P\subseteq[N]$ of modulus $q\leq Q$,
\begin{equation}\label{eq:ga-wtrick-ideal}
 \left|\sum_{n\in P}T_\chi(Wn+b)F_0(g(n)\Gamma)\right|
 \ll_{H_0,W,m,d}QN\log(2N)^7\delta^{ac_E}
                     \|F_0\|_{\mathrm{Lip}}.
\end{equation}
The same implication holds for $E=\mathbb Q$ and the corresponding
ordinary von Mangoldt estimate.
\end{lemma}
\begin{proof}
Take $X=WN+b$ and the sequence $\widetilde g$ constructed above. The
conditions $t=Wn+b$ and $n\in P$ specify an interval progression $R$
of modulus $Wq$. Apply \eqref{eq:ga-ideal-input} with
$Q_1=\max\{Q,Wq,2\}\ll_WQ$ and $\eta=\delta^a$. The ideal sum
then equals the left side of \eqref{eq:ga-wtrick-ideal} exactly, while
$X\ll_WN$ and $\log(2X)\ll_W\log(2N)$. No assertion about
fractional-argument distribution has been assumed: it was proved in
Lemma~\ref{lem:ga-affine-interpolation}.
\end{proof}

\begin{corollary}[Centered ray-prime weights]\label{cor:ga-centered-primes}
Suppose a conjugation-invariant fixed Gaussian ray condition has the
checked finite-character expansion
\[
 \Lambda'_C(t)=\frac12\sum_\chi c_\chi T_\chi(t)+e(t),
 \qquad \sum_{t\leq X}|e(t)|\ll_{H_0}\sqrt X\log(2X)^2.
\]
Assume both the Gaussian estimate \eqref{eq:ga-ideal-input} and its
ordinary rational counterpart, with positive exponents $c_E,c_{\mathbb Q}$.
Let $G/\Gamma$ and $g$ satisfy the hypotheses of
Lemma~\ref{lem:ga-affine-interpolation}, and let $F_0$ be Lipschitz with
Haar mean zero.
Then any fixed linear combination of the $W$-tricked ray-prime weight
and the $W$-tricked ordinary-prime weight obeys, for every interval
progression $P\subseteq[N]$ of modulus at most $Q$,
\[
 \left|\sum_{n\in P}w_N(n)F_0(g(n)\Gamma)\right|
 \ll QN\log(2N)^7\delta^{c_W}\|F_0\|_{\mathrm{Lip}},
 \qquad c_W=a\min\{c_E,c_{\mathbb Q}\}>0,
\]
for mean-zero $F_0$, in a regime
$\delta\geq N^{-\sigma_1}$ and $Q\leq\delta^{-a}$ with fixed
$\sigma_1>0$. Constants may depend on $H_0,W$ and the fixed linear
combination. In particular this applies to the proposed centering
\[
 w_N(n)=\frac{\varphi(W)}W\Lambda'_C(Wn+b)
             -\mu_C(b)\lambda_{b,W}(n),
\]
once the normalization of $\lambda_{b,W}$, the scalar $\mu_C(b)$ and
the displayed character expansion have been fixed and checked.
\end{corollary}
\begin{proof}
Apply Lemma~\ref{lem:ga-wtrick-ma} to the finitely many characters and
to the rational von Mangoldt function. The Gaussian factor $1/2$
accounts for the two conjugate prime ideals above an eligible split
rational prime. Replacing ordinary von Mangoldt by its prime-only
version also has a proper-prime-power error
$O_W(\sqrt N\log(2N)^2)$. Thus the stated estimate holds with the
additional error
$O_{H_0,W}(\sqrt N\log(2N)^2\|F_0\|_\infty)$.
Choose $\sigma_1\leq\sigma$ so small that $\sigma_1c_W<1/4$.
Then $N\delta^{c_W}\geq N^{3/4}$ throughout the claimed regime, and
the square-root errors are absorbed. All normalization factors are
fixed constants because $W$ is fixed. This proves the minor-arc
estimate but does not assert the centered progression estimate
\eqref{eq:gf-sw}, which must be proved separately from the arithmetic
main terms.
\end{proof}

\begin{remark}\label{rem:ga-choice-order}
Fix the arithmetic moduli first, then the positive exponents in this
subsection. They do not depend on the later factorization exponent
$B$. Proposition~\ref{prop:gf-factorization} chooses $B$ large enough
to satisfy the resulting upper complexity condition. Its output
$M$ is polylogarithmic in $N$, so its equidistribution parameter also
obeys $\delta\geq N^{-\sigma_1}$ for sufficiently large $N$.
The affine interpolation argument is motivated by the corresponding
step of Frei--Tinkova, but uses no theorem from that work as an input.
\end{remark}


% Section fragment. Bibliography: GTOrbits2012, GTMobius2012.
\subsection{Quantitative factorization for centered weights}
\label{sec:gf-factorization}

We isolate the analytic implication needed after the arithmetic
progression main terms have been established. Throughout,
$L=\log(2N)$, all nilpotent groups are connected and simply connected,
and all filtrations are by closed connected rational subgroups. Constants
may depend on fixed arithmetic moduli and on the fixed dimension and
degree bounds.

\begin{proposition}[Factorization implication]\label{prop:gf-factorization}
Fix dimension and degree bounds $m,d$. Let $w_N:[N]\to\mathbb C$ satisfy
the following three hypotheses.
\begin{enumerate}
\item For a fixed $K_w\geq0$, $|w_N(n)|\ll L^{K_w}$.
\item For every $A_0,C_0>0$, uniformly over all interval progressions
$P\subseteq[N]$ of common difference at most $L^{C_0}$,
\begin{equation}\label{eq:gf-sw}
 \left|\sum_{n\in P}w_N(n)\right|
       \ll_{A_0,C_0}N L^{-A_0}.
\end{equation}
No lower bound on $|P|$ is imposed here.
\item There are fixed $c>0$, $C_1,K_1\geq0$ such that if a nilmanifold
of dimension at most $m$ and filtration degree at most $d$ has a
$Q$-rational adapted Mal'cev basis, $h$ is a polynomial sequence totally
$\delta$-equidistributed on $[N]$, and $F_0$ is Lipschitz with Haar mean
zero, then
\begin{equation}\label{eq:gf-ma}
 \left|\sum_{n\in P}w_N(n)F_0(h(n)\Gamma)\right|
 \ll N\delta^c Q^{C_1}L^{K_1}\|F_0\|_{\mathrm{Lip}}
\end{equation}
for every interval progression $P\subseteq[N]$ of modulus at most $Q$
and length at least $N/Q$. It is sufficient for this hypothesis to hold
only when
\begin{equation}\label{eq:gf-ma-regime}
 \delta\geq N^{-\sigma},\qquad Q\leq\delta^{-\kappa}
\end{equation}
for fixed positive $\sigma,\kappa$. An estimate allowing arbitrary
progression lengths is stronger than required.
\end{enumerate}
Then, for every $A,C_{\mathrm{init}}>0$, uniformly over initial
nilmanifolds of dimension at most $m$ and degree at most $d$ with an
$L^{C_{\mathrm{init}}}$-rational adapted basis, all polynomial sequences
$g$, and all Lipschitz $F$,
\begin{equation}\label{eq:gf-conclusion}
 \left|\sum_{n\leq N}w_N(n)F(g(n)\Gamma)\right|
     \ll_{A,C_{\mathrm{init}}}N L^{-A}\|F\|_{\mathrm{Lip}}.
\end{equation}
\end{proposition}

\begin{proof}
Normalize $\|F\|_{\mathrm{Lip}}=1$ and put
\[
 D=A+K_w+3,\qquad M_0=L^{\max\{1,C_{\mathrm{init}}\}}.
\]
Leave a factorization exponent $B$ unspecified for the moment. By
\cite[Theorem 1.19]{GTOrbits2012}, there is an integer $M$ with
\[
 M_0\leq M\ll_{B,m,d}M_0^{E(B,m,d)}
\]
and a factorization $g=\varepsilon g'\gamma$ into polynomial sequences.
Here $\varepsilon$ is $(M,N)$-smooth, $\gamma(n)\Gamma$ has period
$q\leq M$, and $g'$ is totally $M^{-B}$-equidistributed in a rational
subgroup quotient. The new adapted basis consists of $M$-rational
combinations of the initial one.

For each residue $j$ modulo $q$, choose a reduced rational representative
$\gamma_j$ of $\gamma(j)\Gamma$. Split that residue class into
$\lceil ML^D\rceil$ consecutive blocks of nearly equal cardinality.
For sufficiently large $N$ after $B$ is fixed, the resulting progressions
$P_{j,k}$ satisfy
\begin{equation}\label{eq:gf-cells}
 \#\{P_{j,k}\}\ll M^2L^D,\qquad
 |P_{j,k}|\gg\frac{N}{M^2L^D},\qquad
 \operatorname{diam}(P_{j,k})\ll\frac{N}{ML^D}.
\end{equation}
Their common difference is $q$. Indeed each original residue class has
cardinality comparable to $N/q$, and rounding the balanced block sizes
changes these estimates only by an absolute factor, since all powers of
$M$ are polylogarithmic in $N$ once $B$ is fixed.

Choose $n_{j,k}\in P_{j,k}$ and put $a_{j,k}=\varepsilon(n_{j,k})$.
The right-invariant ambient metric and smoothness imply
\begin{align*}
 &|F(\varepsilon(n)g'(n)\gamma_j\Gamma)
       -F(a_{j,k}g'(n)\gamma_j\Gamma)|\\
 &\hspace{1cm}\leq d(\varepsilon(n),\varepsilon(n_{j,k}))
       \ll L^{-D}\qquad(n\in P_{j,k}).
\end{align*}
There is no factor depending on the coordinates of $g'(n)$ in this
estimate: right invariance removes it. The total weighted freezing
error is therefore $O(NL^{K_w-D})=O(NL^{-A-3})$.

We use the quantitative conjugation claim proved in
\cite[p.~547]{moebius}, whose inputs are the rational-subgroup,
rational-element and metric comparisons of
\cite[Appendix A]{GTOrbits2012}. There are fixed $a>0$, $C_2\geq1$,
depending only on $m,d$, such that, with
\[
 H_j=\gamma_j^{-1}G'\gamma_j,\quad
 \Lambda_j=H_j\cap\Gamma,\quad
 h_j(n)=\gamma_j^{-1}g'(n)\gamma_j,
\]
the quotient $H_j/\Lambda_j$ has an adapted basis of rationality at most
$M^{C_2}$, $h_j$ is totally $M^{-aB+C_2}$-equidistributed, and
\[
 F_{j,k}(x\Lambda_j)=F(a_{j,k}\gamma_jx\Gamma)
 \quad\hbox{satisfies}\quad
 \|F_{j,k}\|_{\mathrm{Lip}}\leq M^{C_2},\qquad
 \|F_{j,k}\|_\infty\leq1.
\]
We have enlarged $C_2$ to absorb fixed multiplicative constants. These
exponents do not depend on $B$. In applying this published claim, one
must account for the commensurable lattice resulting from rational
conjugation; conjugation is not asserted to be an isometry of the
original quotients. Its finite-index and metric losses are included in
$a,C_2$.

Write $F_{j,k}=m_{j,k}+F_{j,k}^0$ by subtracting its Haar mean on
$H_j/\Lambda_j$. Then $|m_{j,k}|\leq1$,
$\int F_{j,k}^0=0$, and
$\|F_{j,k}^0\|_{\mathrm{Lip}}\leq2M^{C_2}$. This subtraction is needed
even if the initial function $F$ happened to have Haar mean zero.
For a dimension-zero factor the centered function simply vanishes.

Put
\[
 U=C_2+D+5,\qquad Q_*=M^U,\qquad b_*=aB-C_2.
\]
Since $M\geq L$, $Q_*$ dominates the rationality, cell modulus and
$N/|P_{j,k}|$. Applying \eqref{eq:gf-ma} to all centered cell functions
and using \eqref{eq:gf-cells} gives a total of at most
\begin{equation}\label{eq:gf-total-ma}
 C N L^{K_1}M^{-b_*c+C_1U+C_2+2+D}.
\end{equation}
Every cell is counted here. In particular, the bound $N$ in the
minor-arc hypothesis has not been replaced by an unsupported cell
cardinality.

We now fix $B$ sufficiently large that
\begin{equation}\label{eq:gf-choice-B}
 b_*>1,\qquad
 b_*c\geq C_1U+C_2+2+D+A+K_1+4,
 \qquad b_*\kappa\geq U+1
\end{equation}
when the restricted regime \eqref{eq:gf-ma-regime} is used. The last
condition is omitted for an unrestricted minor-arc estimate. These
inequalities involve only fixed exponents and not $E(B,m,d)$.
They make \eqref{eq:gf-total-ma} $O(NL^{-A-4})$. They also ensure
$Q_*\leq(M^{-b_*})^{-\kappa}$. The other regime condition,
$M^{-b_*}\geq N^{-\sigma}$, holds for all sufficiently large $N$:
after fixing $B$ the upper bound on $M$ is a fixed power of $\log N$.

Only at this point use that upper bound, writing $M\leq L^{E_*}$ after
enlarging the fixed exponent $E_*$. There are at most
$O(L^{2E_*+D})$ cells of modulus at most $L^{E_*}$. Apply
\eqref{eq:gf-sw} with
\[
 C_0=E_*+1,\qquad A_0=A+2E_*+D+4.
\]
Summing the constant terms $m_{j,k}\sum_{n\in P_{j,k}}w_N(n)$ gives
$O(NL^{-A-4})$. Combined with the centered estimates and freezing error,
this proves \eqref{eq:gf-conclusion}.
\end{proof}

\begin{remark}[Order of parameter choices]\label{rem:gf-order}
The order in the proof is
\[
 (A,K_w,C_{\mathrm{init}},c,C_1,K_1,a,C_2,\kappa,\sigma)
 \longrightarrow (D,U)\longrightarrow B\longrightarrow E_*
 \longrightarrow(A_0,C_0)\longrightarrow N_0.
\]
Thus there is no requirement that $B$ dominate the factorization upper
exponent $E(B,m,d)$. The centered progression estimate allows every
logarithmic saving, so its exponent can be chosen after $E_*$ is known.
The minor-arc input must have a fixed power form such as
\eqref{eq:gf-ma}; unspecified polylogarithmic specializations alone
would not justify this argument.
\end{remark}


\begin{corollary}[Fixed-modulus centered nilsequence estimate]
\label{cor:ga-fixed-centered-nilsequence}
Fix the Gaussian ray field and a conjugation-invariant union $C$ of ray
classes as above. Let $W$ be a fixed multiple of the conductor of
$K^{\mathrm{ab}}/\Q$, and let $b$ be a unit residue modulo $W$.
For every $A,C_{\mathrm{init}}>0$ and fixed dimension and degree bounds,
uniformly over the nilmanifolds and polynomial sequences in
Proposition~\ref{prop:gf-factorization} with an
$(\log(2N))^{C_{\mathrm{init}}}$-rational adapted basis,
\[
 \left|\sum_{n\leq N}h_{C,b,W}(n)F(g(n)\Gamma)\right|
 \ll_{H,W,A,C_{\mathrm{init}},m,d}
 N(\log(2N))^{-A}\|F\|_{\mathrm{Lip}}.
\]
The constants are uniform in $b$ and in the polynomial orbit. No
uniformity as $W$ varies is asserted.
\end{corollary}
\begin{proof}
We check the three hypotheses of
Proposition~\ref{prop:gf-factorization}. The pointwise inequality
$|h_{C,b,W}|\leq\lambda_{b,W}$ gives its first hypothesis with $K_w=1$.
Corollary~\ref{cor:gaussian-centered-progressions} gives its second
hypothesis for every $A_0,C_0$, in fact for every interval progression.

For the third hypothesis, take $H_0=H$ in
Corollary~\ref{cor:ga-centered-primes}. Its character expansion is
\eqref{eq:gaussian-ray-fourier}, with the factor $1/2$ and the
prime-power error supplied by \eqref{eq:gaussian-ray-prime-factor}
and the bound following that equation. Its two assumed minor-arc
estimates are Proposition~\ref{gi-minor-arc} and
Corollary~\ref{gi-ordinary-primes}. Apply that corollary to the fixed
linear combination defining $h_{C,b,W}$ in
\eqref{eq:gaussian-centered-weight}. It gives \eqref{eq:gf-ma} with
$c=c_W$, $C_1=1$ and $K_1=7$ in the regime
$\delta\geq N^{-\sigma_1}$, $Q\leq\delta^{-a}$.
The affine-interpolation result initially restricts $\delta$ to a
fixed sufficiently small interval. For larger $\delta$, the bound
$|h_{C,b,W}|\ll_W\log(2N)$ gives \eqref{eq:gf-ma} directly after
enlarging its fixed constant. For a dimension-zero quotient the
mean-zero test function vanishes. For all other dimensions and degrees
within the fixed bounds, take the minimum of the finitely many positive
exponents $c_W,a,\sigma_1$ and the maximum of the fixed constants.
Thus the restricted regime \eqref{eq:gf-ma-regime} holds with fixed
positive $\kappa,\sigma$, uniformly throughout the required class.

Proposition~\ref{prop:gf-factorization} now proves the displayed
estimate. There are only finitely many unit residues modulo the fixed
$W$, so their constants can be bounded by a single constant. In
particular, division by $N$ and the relabeling $W=V$, $N=L$ give the
fixed-complexity estimate \eqref{eq:gaussian-centered-nil-input}.
\end{proof}


% The analytic bridge below uses the fixed-modulus centered nilsequence
% estimate just proved for the Gaussian ray fields under consideration.
